\chapter{Where the State Model Ends and Dynamics Begins}
\label{chap:concept-boundaries}

Part V already established pose2, pose3, moving frames, cant, Metric
Realization, target and physical state, and observation. This chapter does not
define them again. It fixes the narrower responsibilities of Part VI so that
later chapters consume one railway-state model.

\section{Inherited Railway State}

The input is the qualified Part V state at \(s\), including its Alignment
reference, realization, frame convention, provenance, applicability, and
validity. ``Track state'' in the following chapters is a short reader-facing
name for the applicable subset of that structure, not a competing state type.
Likewise, a represented coordinate tuple is not a replacement for pose3 and an
observed pose is not a physical state.

\section{New Responsibilities}

\begin{description}
\item[Operational velocity] qualifies the traversal law \(s(t)\) and its
derivatives.
\item[Temporal exposure] composes spatial curvature, cant, profile, and frame
quantities with that traversal.
\item[Reduced kinematics] evaluates assumption-bound indicators without
claiming complete response.
\item[Vehicle space] expresses vehicle-related quantities relative to the
qualified track frame and a stated vehicle model.
\item[Response modelling] adds the physical parameters needed for a predicted
vehicle--track response.
\item[Dynamic balance] applies a declared reduced relation under its stated
frame, sign, speed, and cant conventions.
\item[Realization-aware dynamics] compares qualified target, realized, or
observed inputs without merging them.
\item[Optimization] varies declared degrees of freedom against explicit
objectives and constraints; it neither supplies objectives nor decides.
\end{description}

\section{One State, Several Uses}

The same qualified railway state may feed a clearance calculation, a reduced
lateral indicator, a multibody model, or an optimization experiment. These
uses produce different outputs because they add different models and purposes.
They do not create alternative definitions of the Alignment or railway state.

\begin{principle}[Single-definition principle]
Pose2, pose3, the moving frame, cant, realization, physical state, and
observation retain their established meanings. Part VI adds traversal,
response models, and purpose-qualified uses by reference.
\end{principle}

\begin{summary}
Part VI begins at \(s(t)\). Everything spatial is inherited; everything
temporal or response-related must identify the additional input that creates
it.
\end{summary}
